\documentclass[12pt, letterpaper]{article} 

\usepackage{amsmath, amssymb, amsthm, enumerate}

\title{MATH 174 - HW 1}
\author{Dani Nguyen}
\date{Fall 2025}

\begin{document}
\maketitle

\begin{enumerate}
    \item %Problem 1
        \begin{enumerate}[(a)]
            \item Forward substitution on $Ly=b$
                \begin{align*}
                    x_1               &= 2  \\
                    -x_1 + x_2        &= -1 \\
                    3x_1 - 3x_2 + x_3 &= 3
                \end{align*}
                So, we have 
                \begin{align*}
                    x_1 &= 2 \\
                    x_2 &= -1 + x_1 \\
                        &= -1 + 2 \\
                        &= 1 \\
                    x_3 &= 3 - 3x_1 + 3x_2 \\
                        &= 3 - 3(2) + 3(1) \\
                        &= 0
                \end{align*}
            \item Backward substitution on $Ux=y$
                \begin{align*}
                    -1x_1 + 2x_2 + x_3 &= y_1  \\
                    2x_2 - 2x_3 &= y_2 \\
                    5x_3 &= y_3
                \end{align*}
                So, we have
                \begin{align*}
                    x_3 &= \frac{y_3}{5} \\
                    2x_2 &= y_2 + 2x_3 \\
                    2x_2 &= y_2 + 2\left(\frac{y_3}{5}\right) \\
                    x_2  &= \frac{1}{2}y_2 + \frac{y_3}{5} \\
                    -x_1 &= y_1 - 2x_2 - x_3 \\
                    x_1 &= -y_1 + 2y_2 + \frac{2y_3}{5} + \frac{y_3}{5} \\
                    &= -y_1 + 2y_2 + \frac{3y_3}{5}
                \end{align*}
            \item Compute $A=LU$
                \[
                    A =
                    \begin{bmatrix}
                        1 & 0 & 0 \\
                        -1 & 1 & 0 \\
                        3 & -3 & 1
                    \end{bmatrix}
                    \begin{bmatrix}
                        -1 & 2 & 3 \\
                        0 & 2 & -2 \\
                        0 & 0 & 5
                    \end{bmatrix}
                    = 
                    \begin{bmatrix}
                        -1 & 2 & 3 \\
                        1 & 0 & -5 \\
                        -3 & 0 & 20
                    \end{bmatrix}
                \]
                \begin{align*}
                    Ax=b : 
                    \begin{bmatrix}
                        -1 & 2 & 3 \\
                        1 & 0 & -5 \\
                        -3 & 0 & 20
                    \end{bmatrix}
                    \begin{bmatrix}
                        x_1 \\
                        x_2 \\
                        x_3
                    \end{bmatrix}
                    &=
                    \begin{bmatrix}
                        2 \\
                        -1 \\
                        3
                    \end{bmatrix} \\
                    \begin{bmatrix}
                        1 & 0 & 0 \\
                        0 & 1 & 0 \\
                        0 & 0 & 1
                    \end{bmatrix}
                    \begin{bmatrix}
                        x_1 \\
                        x_2 \\
                        x_3
                    \end{bmatrix}
                    &=
                    \begin{bmatrix}
                        -1 \\
                        \frac{1}{2} \\
                        0
                    \end{bmatrix}
                \end{align*}
                The solution to $Ax+b$ is any multiple of $\begin{bmatrix} -1 \\ \frac{1}{2} \\ 0 \end{bmatrix}$
        \end{enumerate}
    \item %Problem 2
        \begin{enumerate}[(a)]
            \item The summation is from $j=1$ to $i-1$, there are $i-1$ additions. There is also one more subtraction between $b_i$ and the summation. So in total, there are $i$ additions and subtractions.
            \item There are $i-1$ terms in the summation, therefore there are $i-1$ multiplications.
        \end{enumerate}
    \item %Problem 3 
        \begin{enumerate}[(a)]  
            \item Gaussian elimination
                \begin{align*}
                    \begin{bmatrix}
                        2  & -3 & 1  & \bigm| & 1\\
                        1  & 1  & -1  & \bigm| & 2\\
                        -4 & 0  & 4 & \bigm| & -1
                    \end{bmatrix}& \\
                    \begin{bmatrix}
                        2 & -3 & 1  & \bigm| & 1\\
                        1 &  1 & -1 & \bigm| & 2\\
                        0 & -6 & 6  & \bigm| & 1
                    \end{bmatrix}& R_3 = R_3 + 2R_1\\
                    \begin{bmatrix}
                        1 & -4 & 2  & \bigm| & -1\\
                        1 & 1  & -1 & \bigm| & 2\\
                        0 & -6 & 6  & \bigm| & 1
                    \end{bmatrix}& R_1 = R_1 - R_2\\
                    \begin{bmatrix}
                        1 & -4 & 2  & \bigm| & -1\\
                        0 & 5  & -3 & \bigm| & 3 \\
                        0 & -6 & 6  & \bigm| & 1
                    \end{bmatrix}& R_2 = R_2 - R_1\\
                    \begin{bmatrix}
                        1 & -4 & 2  & \bigm| & -1\\
                        0 & -1 & 3  & \bigm| & 4 \\
                        0 & -6 & 6  & \bigm| & 1
                    \end{bmatrix}& R_2 = R_2 + R_3\\
                    \begin{bmatrix}
                        1 & -4 & 2  & \bigm| & -1\\
                        0 &  1 & -3  & \bigm| & -4 \\
                        0 & -6 & 6  & \bigm| & 1
                    \end{bmatrix}& R_2 = -R_2 \\
                    \begin{bmatrix}
                        1 & -4 & 2  & \bigm| & -1\\
                        0 &  1 & -3  & \bigm| & -4 \\
                        0 & 0 & -12  & \bigm| & -23
                    \end{bmatrix}& R_3 = R_3+6R_2 \\
                    \begin{bmatrix}
                        1 & -4 & 2  & \bigm| & -1\\
                        0 &  1 & -3  & \bigm| & -4 \\
                        0 & 0 & 1  & \bigm| & \frac{23}{12}
                    \end{bmatrix}& R_3 = \frac{-1}{12}R_3 \\
                \end{align*}
                Back substitution
                \begin{align*}
                    x_3 &= \frac{23}{12} \\
                    x_2 &= 3x_3 - 4 \\
                        &= \frac{23}{4}-4 \\
                        &= \frac{7}{4} \\
                    x_1 &= 4x_2 - 2x_3 - 1 \\
                        &= 7 - \frac{23}{6} - 1 \\ 
                        &= \frac{13}{6}
                \end{align*}
        \end{enumerate}
    \item %Problem 4
        \begin{enumerate}[(a)]
            \item Gaussian elimination
                \begin{align*}
                    \begin{bmatrix}
                        1     &   4 &  \bigm| & 3\\
                        0.002 &  -4 &  \bigm| & -2 
                    \end{bmatrix}& \\
                    \begin{bmatrix}
                        1 & 4 &  \bigm| & 3\\
                        0 & 4 &  \bigm| & -1\times 10^3 
                    \end{bmatrix}& \\
                \end{align*}
            \item 
            \item 
        \end{enumerate}
                
    \item %Problem 5
    \item %Problem 6
    \item %Problem 7
    \item %Problem 8
\end{enumerate}
\end{document}